\title{Direct Preference Optimization: Your Language Model is Secretly a Reward Model}

\begin{abstract}
Large-scale unsupervised language models (LMs) are proficient at acquiring extensive world knowledge and certain reasoning capabilities, yet exerting fine-grained control over their outputs is challenging due to their entirely unsupervised training regime. Current approaches to enhance this controllability typically involve gathering human judgments about the relative merits of generated texts and subsequently adapting the LMs through fine-tuning, commonly via reinforcement learning from human feedback (RLHF). RLHF, however, is intricate and prone to instability: it first constructs a reward model that encapsulates human preferences, and then applies reinforcement learning to adjust the unsupervised LM so as to optimize this modeled reward—while constraining deviation from the pretrained model. This work proposes an alternative parameterization of the reward model in RLHF, which facilitates direct, closed-form derivation of the optimal policy. As a result, the canonical RLHF objective can be addressed using a straightforward classification loss. The method, named \textit{Direct Preference Optimization} (DPO), offers a robust, efficient, and lightweight alternative, circumventing the necessity for LM sampling during adaptation or extensive tuning of hyperparameters. Empirical evaluations demonstrate that DPO is at least as effective as prevalent techniques for aligning LMs with human preferences. Notably, DPO surpasses RLHF approaches based on PPO in steering generation sentiment, and achieves comparable or superior response quality in tasks such as summarization and single-turn conversation, all while remaining markedly easier to implement and train.
\end{abstract}

\section{Introduction}
Extensive unsupervised training of large language models (LMs) on massive corpora leads to the emergence of unexpected proficiencies~\citep{chowdhery2022palm, brown2020language, touvron2023llama,bubeck2023sparks}. However, because these corpora are sourced from human-authored content that reflects a broad spectrum of intentions, priorities, and expertise, not all such content is ideal for replication by language models. For example, while an AI coding assistant should be capable of \textit{recognizing} typical programming errors in order to offer corrections, it is preferable for its code generation to be informed by the infrequent but high-caliber examples within its training dataset. Likewise, it is beneficial for a language model to be \textit{knowledgeable} about widely held misconceptions, but it should avoid echoing such falsehoods proportionally to their prevalence in its training data. Thus, the task of aligning the model's \emph{behavior and outputs} with desirable qualities drawn from its extensive \textit{latent knowledge and skill set} is essential for developing safe, high-performing, and reliably controlled AI systems~\citep{ouyang2022training}. Traditionally, techniques for shaping language model behavior are based on reinforcement learning (RL) to match human preferences. In this work, we demonstrate that the objectives pursued by RL-based alignment techniques can, in fact, be optimized exactly through a straightforward binary cross-entropy loss, significantly streamlining the preference learning process.

Broadly, conventional approaches introduce desired behavioral tendencies into language models via collections of human-annotated preferences that encode criteria for safety and usefulness. This phase of preference learning is performed following a preliminary, unsupervised pre-training phase on vast textual data. Although supervised fine-tuning with high-quality human demonstrations represents a simple pathway for preference alignment, reinforcement learning from human or AI-generated feedback (RLHF/RLAIF; \citep{christiano2017deep,bai2022constitutional}) has proven most effective. RLHF constructs a reward model trained on datasets of human preferences and leverages RL to adjust the model’s policy to yield responses judged as high-reward, all while maintaining proximity to the original model’s behavior. Despite RLHF's success in cultivating advanced conversational and programming skills, its implementation remains considerably more intricate than supervised fine-tuning, necessitating the training of multiple LMs and policy sampling throughout the process, resulting in elevated computational demands.

This work demonstrates a method for directly training language models to align with human preferences, circumventing the need for explicit reward modeling or reinforcement learning techniques. We introduce \textit{Direct Preference Optimization} (DPO), a novel algorithm that, while targeting the same reward maximization objective subject to a KL-divergence constraint as conventional RLHF approaches, offers a more accessible implementation and streamlined training process. Conceptually, DPO modifies the log probability ratio between preferred and non-preferred outputs, but introduces a dynamic, example-specific importance weighting mechanism that mitigates the model degradation observed with straightforward probability ratio objectives. DPO, much like earlier algorithms, utilizes a theoretical preference framework—such as the Bradley-Terry model \cite{bradley1952rankanalysis}—to evaluate how accurately a reward function mirrors observed human preference data. Unlike previous pipelines that employ the preference model to define a preference loss for reward model training before policy optimization, DPO leverages a change of variables, allowing the preference loss to depend directly on the policy parameters. With access to a dataset encoding human choices among candidate outputs, DPO optimizes the model using a straightforward binary cross entropy loss, yielding a policy optimal with respect to an implicitly learned reward function that best fits the preference observations.

Our primary contribution is Direct Preference Optimization (DPO): a straightforward algorithm, free from reinforcement learning, designed to train language models based on preference data. Through experiments, we demonstrate that DPO matches or surpasses the efficacy of established approaches, including PPO-based RLHF, in applications such as sentiment modification, summarization, and conversational tasks, using models with up to 6B parameters.

\section{Related Work}

Large self-supervised language models have demonstrated the capacity to address certain tasks in a zero-shot fashion \citep{radford2019language} or through few-shot prompting \citep{gpt3,megatron,chowdhery2022palm}. Nevertheless, adapting these models to downstream tasks and ensuring their responses better reflect user intentions can be substantially enhanced by further fine-tuning on curated instruction datasets and corresponding human-generated completions \citep{mishra-etal-2022-cross,sanh2022multitask,chung2022scaling,thoppilan2022lamda}. This approach, known as `instruction-tuning', broadens the generalization abilities of LLMs to instructions outside the original training set, thereby improving overall usability \citep{chung2022scaling}. While instruction tuning is effective, collecting \textit{relative} human evaluations of output quality is often less resource-intensive than obtaining expert-labeled examples. Consequently, recent research has focused on adapting LLMs using preference datasets derived from human judgments, yielding gains in tasks like translation \citep{kreutzer-etal-2018-reliability}, summarization \citep{stiennon2022learning,ziegler2020finetuning}, narrative generation \citep{ziegler2020finetuning}, and instruction compliance \citep{ouyang2022training,ramamurthy2023is}. In these frameworks, the first step involves training a neural reward model to align with collected preference data under models like Bradley-Terry \citep{bradley1952rankanalysis}, followed by optimizing the language model with reinforcement learning algorithms—typically REINFORCE \citep{williams1992reinforce}, proximal policy optimization (PPO; \cite{schulman2017proximal}), or similar strategies \citep{ramamurthy2023is}—to maximize reward. A related research direction uses instruction-following LLMs, fine-tuned with human feedback, to create new synthetic preference datasets for specific criteria such as safety or harmlessness \citep{bai2022constitutional}, often relying on human-authored text rubrics to guide the model's annotations through weak supervision. Collectively, these strategies draw from two main research traditions: reinforcement learning for language model training across diverse objectives~\citep{Ranzato2015SequenceLT,paulus2018a,wu2018learning}, and general frameworks for learning from human preference signals \citep{christiano2017deep,kupcsik2018learning}. Yet, despite the practical attractiveness of utilizing comparative human preferences, reinforcement learning-based fine-tuning of large language models continues to pose significant logistical difficulties. The present work addresses this by proposing a theoretically sound alternative for optimizing models with relative preferences, circumventing the need for RL.

Beyond applications in language, the task of learning policies from preferences has been examined within both bandit and reinforcement learning frameworks, leading to the development of various methodologies. When preferences or rankings over actions, rather than explicit rewards, are employed in contextual bandit learning, this scenario is referred to as the contextual dueling bandit (CDB; \cite{yue2012karmed,dudik2015contextual}). In cases where absolute reward values are unavailable, theoretical investigations of CDBs replace the concept of an optimal policy with that of a \textit{von Neumann winner}: a policy that, on average, achieves at least a 50\% win rate against \textit{any} competing policy \citep{dudik2015contextual}. Notably, CDB frameworks involve receiving preference annotations in an online manner, whereas the human preference learning paradigm typically utilizes a fixed collection of action pairs annotated offline with preference labels \citep{yan2022human}. Likewise, \textit{preference-based RL} (PbRL) addresses learning from binary preferences that originate from an \textit{unknown} `scoring' function instead of direct rewards \citep{BusaFekete2014,ruiz2023dueling}. PbRL offers several algorithmic solutions, with some leveraging off-policy preference data; however, these approaches commonly entail the explicit inference of the underlying scoring (or reward) function, followed by a policy optimization step based on this estimate \citep{jain2013learning,BusaFekete2014,christiano2017deep,sadigh2017active,kupcsik2018learning}. In contrast, our work proposes a unified policy learning method that directly trains policies to comply with given preferences in a single stage.

\section{Preliminaries}\label{section:prelims}

We provide an overview of the RLHF pipeline as described in \citeauthor{ziegler2020finetuning} (also referenced in \citep{stiennon2022learning, bai2022training, ouyang2022training}), which conventionally comprises three sequential components: (1) supervised fine-tuning (SFT); (2) acquisition of preference data and learning of the reward model; and (3) policy optimization via reinforcement learning.

\textbf{SFT}: RLHF workflows customarily initiate with the fine-tuning of a pre-trained language model by means of supervised learning on high-quality examples relevant to the target application (such as dialogue or summarization), resulting in a model $\pi^\text{SFT}$.

\textbf{Reward Modelling Phase}: In the subsequent phase, prompts $x$ are given to the SFT model, which generates pairs of outputs $(y_1, y_2)\sim \pi^\text{SFT}(y \mid x)$. These paired responses are evaluated by human annotators, who indicate a preference for one response, formalized as $y_w\succ y_l \mid x$ where $y_w$ represents the chosen (preferred) answer and $y_l$ the alternative, less favored response within $(y_1, y_2)$. The assumption is that these preferences reflect an underlying, unobserved reward function $r^*(y, x)$. Various models are employed for representing such preference data, with the Bradley-Terry (BT) model \cite{bradley1952rankanalysis} being a common selection (although the more flexible Plackett-Luce model \citep{plackett1975analysis, luce2012individual} is also applicable when ranked sets beyond pairs are available). In the BT framework, the human preference probability distribution $p^*$ can be formalized as:

\begin{equation}\label{eq:bradley-terry}
    p^*(y_1\succ y_2 \mid x)=\frac{\exp\left(r^*(x, y_1)\right)}{\exp\left(r^*(x, y_1)\right) + \exp\left(r^*(x, y_2)\right)}.
\end{equation}

Given a fixed dataset of pairwise comparisons $\mathcal{D}=\bigl\{x^{(i)}, y_w^{(i)}, y_l^{(i)}\bigr\}_{i=1}^N$ sampled according to $p^*$, we can introduce a parametric reward model $r_{\phi}(x, y)$ and fit its parameters by maximizing the likelihood. If we interpret the setup as a binary classification task, the corresponding negative log-likelihood objective is:

\begin{equation}\label{eq:reward_model}
    \mathcal{L}_R(r_{\phi}, \mathcal{D}) = -\mathbb{E}_{(x, y_w, y_l)\sim \mathcal{D}}\bigl[\log \sigma(r_{\phi}(x, y_w)- r_{\phi}(x, y_l))\bigr]
\end{equation}

where $\sigma$ denotes the logistic sigmoid function. In the context of language models, the reward estimator $r_{\phi}(x, y)$ is generally initialized from the SFT model $\pi^\text{SFT}(y \mid x)$, commonly by appending a linear projection to the final transformer layer that outputs a scalar reward estimate \cite{ziegler2020finetuning}. To reduce reward variance, prior studies apply normalization such that $\mathbb{E}_{x,y\sim \mathcal{D}}\left[r_\phi(x, y)\right] = 0$ for all $x$.

\textbf{RL Fine-Tuning Phase}: In the RL fine-tuning stage, the learned reward function delivers evaluative signals to guide language model outputs. Following established work~\citep{jaques2017sequence, jaques2020human}, the optimization objective takes the form:

\begin{equation}\label{eq:RL}
\max_{\pi_{\theta}}  \mathbb{E}_{x\sim \mathcal{D}, y\sim \pi_{\theta}(y \mid x)}\bigl[r_{\phi}(x, y)\bigr] - \beta\mathbb{D}_{\textrm{KL}}\bigl[\pi_{\theta}(y\mid x)\mid \mid \pi_\text{ref}(y\mid x)\bigr],
\end{equation}

where $\beta$ tunes the strength of the regularization term relative to a reference policy $\pi_\text{ref}$—typically chosen as the initial SFT model $\pi^\text{SFT}$. The policy $\pi_\theta$ is also usually initialized from $\pi^\text{SFT}$. This regularization term is crucial to prevent the policy from straying outside the distribution where the reward model produces reliable estimates, thereby preserving output diversity and mitigating mode collapse toward isolated, high-reward responses. Because language generation is discrete and non-differentiable, reinforcement learning methods are typically employed to optimize this objective. Standard practice \citep{ziegler2020finetuning, stiennon2022learning, bai2022training, ouyang2022training} is to define the reward as ${r(x, y) = r_{\phi}(x, y) -\beta (\log \pi_{\theta}(y\mid x) - \log \pi_\text{ref}(y\mid x))}$, which is then optimized using PPO \cite{schulman2017proximal}.

\section{Direct Preference Optimization}\label{sec:DPO}

Given the complexity and resource requirements associated with scaling reinforcement learning techniques for language model fine-tuning, our objective is to establish a more direct and tractable framework for optimizing policies with human preference data. Contrary to standard RLHF pipelines—which require separate reward learning and subsequent RL-based policy optimization—our approach is predicated on a carefully chosen reward function parameterization, which makes it possible to recover the optimal policy in closed form, sidestepping the need for a separate reinforcement learning stage. As detailed below, the central insight is to exploit an analytic correspondence between reward models and their induced optimal policies, thereby recasting optimization over rewards as direct optimization over policies. By adopting this change-of-variables perspective, we can forgo explicit reward model estimation while still adhering to well-established human preference models, such as the Bradley-Terry framework. Under this method, the policy network encodes both the generative

\textbf{Deriving the DPO objective.} We begin from the same reinforcement learning objective as in prior studies, Eq.~\ref{eq:RL}, and assume a general reward function $r$. Prior research~\citep{peters2007reinforcement, peng2019advantage, korbak2022reinforcement, go2023aligning} has shown that the solution to the KL-regularized reward maximization problem in Eq.~\ref{eq:RL} has the form:

\begin{equation}\label{eq:op_policy}
    \pi_r(y\mid x) = \frac{1}{Z(x)}\pi_\text{ref}(y\mid x)\exp\left(\frac{1}{\beta}r(x, y)\right),
\end{equation}

Here, $Z(x) =\sum_{y}\pi_\text{ref}(y\mid x)\exp\left(\frac{1}{\beta}r(x, y)\right)$ denotes the partition function. A complete proof can be found in Appendix \ref{app:derivation1}. Even if the reward function $r^*$ is approximated by an MLE estimate $r_{\phi}$, direct computation of the partition function $Z(x)$ remains challenging~\citep{korbak2022reinforcement, go2023aligning}, which limits practical usability of this formulation. Nonetheless, Eq.~\ref{eq:op_policy} can be manipulated to express $r$ as a function of its optimal policy $\pi_r$, the reference $\pi_\text{ref}$, and $Z(\cdot)$. By taking the logarithm of both sides of Eq.~\ref{eq:op_policy}, we have:

\begin{equation}\label{eq:main_eq}
    r(x,y) =\beta \log \frac{\pi_r(y\mid x)}{\pi_\text{ref}(y\mid x)} + \beta \log Z(x).
\end{equation}

This reparametrization applies equally to the ground-truth reward $r^*$ and the associated optimal policy $\pi^*$. The Bradley-Terry model is determined by reward differences between two outputs: ${p^*(y_1 \succ y_2 \mid x) = \sigma(r^*(x, y_1) - r^*(x, y_2))}$. Plugging in the form from Eq.~\ref{eq:main_eq} into the preference likelihood, as defined in Eq.~\ref{eq:bradley-terry}, the terms involving the partition function $Z(x)$ cancel. As a result, the preference probability depends solely on the optimal policy $\pi^*$ and reference $\pi_\text{ref}$. This gives that the RLHF-optimal policy $\pi^*$ satisfies:

\begin{equation}\label{eq:objective}
    p^*(y_1\succ y_2 \mid x)=\frac{1}{1 + \exp\left(\beta \log \frac{\pi^*(y_2\mid x)}{\pi_\text{ref}(y_2\mid x)} - \beta \log \frac{\pi^*(y_1\mid x)}{\pi_\text{ref}(y_1\mid x)}\right)}
\end{equation}

Further steps can be found in Appendix~\ref{app:derivation2}. While Eq.~\ref{eq:objective} assumes the Bradley-Terry model, analogous derivations exist for the more general Plackett-Luce models~\citep{plackett1975analysis, luce2012individual}, as elaborated in Appendix~\ref{app:plackett_luce_models}.

Given this characterization of preference probabilities in terms of the optimal policy, rather than the underlying reward, we can construct a maximum likelihood estimation problem over a parameterized policy $\pi_\theta$. Paralleling the approach used in reward modeling (see Eq.~\ref{eq:reward_model}), the policy-based objective is thus:

\begin{equation}\label{eq:optimum_model}
    \mathcal{L}_\text{DPO}(\pi_{\theta}; \pi_\text{ref}) = -\mathbb{E}_{(x, y_w, y_l)\sim \mathcal{D}}\left[\log \sigma \left(\beta \log \frac{\pi_{\theta}(y_w\mid x)}{\pi_\text{ref}(y_w\mid x)} - \beta \log \frac{\pi_{\theta}(y_l\mid x)}{\pi_\text{ref}(y_l\mid x)}\right)\right].
\end{equation}

By taking this approach, we approximate an implicit reward via an alternative parameterization, for which the optimal policy corresponds directly to $\pi_\theta$. Furthermore, because this method is mathematically equivalent to fitting a reparametrized Bradley-Terry model, it inherits several beneficial theoretical attributes, including consistency under certain assumptions regarding the distribution of preference data \cite{bong2022generalized}. In Section~\ref{sec:theory}, we elaborate on the theoretical underpinnings of DPO and how they relate to prior literature.

\textbf{Mechanism of the DPO Update.} To gain insight into the functioning of DPO, it is instructive to examine the gradient of its objective function, $\mathcal{L}_\text{DPO}$. The parameter gradient can be expressed as:

\begin{multline*}\label{eq:gradient}
    \nabla_\theta \mathcal{L}_\text{DPO}(\pi_\theta;\pi_\text{ref}) = \\ -\beta\mathbb{E}_{(x, y_w, y_l) \sim \mathcal{D}} \bigg[\underbrace{\sigma(\hat{r}_\theta(x, y_l) - \hat{r}_\theta (x, y_w))}_\text{greater effect when reward ranking is inaccurate}\bigg[\underbrace{\nabla_\theta\log \pi(y_w \mid x)}_\text{promote $y_w$ probability} - \underbrace{\nabla_\theta\log\pi(y_l \mid x)}_\text{reduce $y_l$ probability}\bigg]\bigg],
\end{multline*}

where $\hat{r}_\theta(x, y) = \beta \log \frac{\pi_\theta(y \mid x)}{\pi_\text{ref}(y \mid x)}$ represents the reward implicitly defined in terms of both the policy $\pi_\theta$ and the reference model $\pi_\text{ref}$ (additional details in Section~\ref{sec:theory}). In essence, this gradient steers $\pi_\theta$ towards assigning greater likelihood to the preferred answers $y_w$, while lowering the likelihood assigned to the less preferred alternatives $y_l$. Crucially, each training sample is weighted by the extent to which the current reward estimator $\hat{r}_\theta$ incorrectly assigns higher reward to the disfavored $y_l$, modulated by the scaling parameter $\beta$. This reflects both the degree of ranking error and the influence of the KL constraint. Empirically, our findings underscore the necessity of this weighting: omitting the weighting factor (i.e., using a na\"ive implementation) can lead to severe degradation of language model quality, as shown in Appendix Table~\ref{tab:unlikelihood_generations}.

\textbf{DPO overview.} The standard DPO procedure comprises the following steps: 1) For each prompt $x$, generate completion samples $y_1, y_2 \sim \pi_\text{ref}(\cdot \mid x)$, and annotate them according to human preference to form the offline preference dataset $\mathcal{D} = \{x^{(i)}, y_w^{(i)}, y_l^{(i)}\}_{i=1}^N$; 2) Update the language model $\pi_\theta$ by minimizing $\mathcal{L}_\text{DPO}$ with respect to the chosen $\pi_\text{ref}$, the constructed dataset $\mathcal{D}$, and a specified $\beta$ parameter. 
In real-world scenarios, practitioners often prefer to leverage publicly available preference datasets instead of producing new samples and collecting human feedback. Because such datasets are typically collected using $\pi^\text{SFT}$, we set $\pi_\text{ref} = \pi^\text{SFT}$ wherever possible. In cases where $\pi^\text{SFT}$ is inaccessible, we determine $\pi_\text{ref}$ by maximizing the likelihood over preferred completions ${(x, y_w)}$, i.e., by solving ${\pi_\text{ref} = \argmax_{\pi}\mathbb{E}_{x, y_w \sim \mathcal{D}}\left[\log \pi(y_w \mid x)\right]}$. This strategy addresses the mismatch between the unavailable true reference distribution and the practical $\pi_\text{ref}$ adopted in DPO. Additional implementation specifics and hyperparameter selections are detailed in Appendix~\ref{app:implementation}.

\section{Theoretical Analysis of DPO} In this part, we offer a deeper theoretical interpretation of the DPO framework, articulate its foundational properties, and clarify how DPO overcomes some limitations observed in actor-critic approaches to RLHF (such as PPO~\cite{schulman2017proximal}).

\label{sec:theory}

\subsection{Your Language Model Is Secretly a Reward Model} DPO eliminates the need to explicitly learn a reward function and perform reinforcement learning to derive the policy, achieving both via a single maximum likelihood objective. Observe that the loss defined in Eq. \ref{eq:main_eq} matches the Bradley-Terry formulation with rewards parameterized as $r^*(x, y) = \beta \log\frac{\pi^*_\theta(y \mid x)}{\pi_\text{ref}(y \mid x)}$, and thus, by optimizing the model $\pi_{\theta}$, we perform a task equivalent to reward model learning as formulated in Eq. \ref{eq:reward_model} under an appropriate variable transformation. In this discussion, we develop the theoretical justification for this parameterization, demonstrate that it retains the expressivity of possible reward models, and show it enables exact recovery of the optimal policy. Our exposition starts by introducing an equivalence relation for reward functions.

\begin{definition}
Two reward functions $r(x, y)$ and $r'(x, y)$ are defined to be equivalent if and only if there exists a function $f$ such that ${r(x, y)-r'(x, y) = f(x)}$.     
\end{definition}

It can be readily verified that this relation satisfies the properties of an equivalence relation and thus divides the space of reward functions into disjoint equivalence classes. We present two key lemmas:

\begin{lemma}\label{lemma:same_prefrence} Within the Plackett-Luce framework—and specifically in the case of the Bradley-Terry model—reward functions belonging to the same equivalence class lead to identical preference distributions.
\end{lemma}

\begin{lemma}\label{lemma:same_policy}
    Any pair of reward functions within the same equivalence class produces the same optimal policy for the constrained reinforcement learning problem.
\end{lemma}

The arguments supporting these statements are elementary and have been deferred to Appendix \ref{app:lemma1}. The first lemma highlights a commonly recognized under-determination inherent to models in the Plackett-Luce family \cite{plackett1975analysis}. As a result of this ambiguity, it is standard practice to enforce additional constraints to ensure identifiability and thereby guarantee meaningful maximum likelihood estimates from Eq. \ref{eq:reward_model} \cite{bong2022generalized}. The second lemma asserts that the entire equivalence class of reward functions yields the same solution for the optimal policy. Accordingly, our focus can be restricted to recovering any representative reward function from the relevant equivalence class. The following theorem, whose proof appears in Appendix~\ref{app:thm1}, formalizes this insight:

\begin{theorem}\label{thm:main}
    Given mild regularity conditions, every equivalence class of rewards consistent with the Plackett-Luce (and, in particular, the Bradley-Terry) model admits a representation of the form ${r(x, y) = \beta \log \frac{\pi(y\mid x)}{\pi_\text{ref}(y\mid x)}}$ for an appropriate model $\pi(y\mid x)$ and a fixed reference model $\pi_\text{ref}(y \mid x)$.
\end{theorem}

\begin{sproof}
    Consider an arbitrary reward function $r(x, y)$, which induces an optimal distribution $\pi_r(y \mid x)$ as given by Eq. \ref{eq:op_policy}. We will establish that there exists a function within the equivalence class of $r$ expressible via the aforementioned reparameterization. Define the transformation $f$ as  
\begin{equation}
    f(r; \pi_\text{ref}, \beta)(x, y) = r(x, y) - \beta\log\sum_{y}\pi_\text{ref}(y\mid x)\exp\left(\frac{1}{\beta}r(x, y)\right)
\end{equation}
This operator $f$ acts to re-center the reward by subtracting the log-partition function corresponding to $\pi_r$ and reference distribution $\pi_\text{ref}$. Because this normalization only involves $x$, the output $f(r; \pi_\text{ref}, \beta)(x, y)$ remains in the equivalence class of $r(x, y)$. Substituting $r$ with the right side of Eq.~\ref{eq:main_eq}, valid for all reward functions, yields $f(r; \pi_\text{ref}, \beta)(x, y) = \beta \log \frac{\pi_r(y\mid x)}{\pi_\text{ref}(y\mid x)}$. Therefore, this projection selects an equivalent reward function of the specified form, establishing that this reparameterization incurs no loss of generality in our model.
\end{sproof}

Theorem~\ref{thm:main} can be reinterpreted as identifying the precise reward function selected by the DPO reparameterization within each equivalence class—specifically, the reward function that fulfills the following condition:

\begin{equation}\label{eq:lag_p}
     \sum_{y}\underbrace{\pi_\text{ref}(y\mid x)\exp\left(\frac{1}{\beta}r(x, y)\right)}_{=\pi(y\mid x)\text{, by Thm.~\ref{thm:main} reparam.}} = 1,
\end{equation}

That is, $\pi(y\mid x)$ forms a proper probability distribution: the probabilities remain non-negative and total to one. Moreover, according to Eq.~\ref{eq:op_policy}, Eq.~\ref{eq:lag_p} serves as the partition function for the optimal policy that arises from the reward function $r(x, y)$. The fundamental insight underlying the DPO algorithm is its ability to enforce particular constraints on the otherwise under-specified Plackett-Luce class of preference models (including Bradley-Terry as a special case). In doing so, the algorithm maintains the representational scope of possible reward models, while ensuring the optimal policy in Eq.~\ref{eq:op_policy} is analytically manageable for any prompt $x$.

\subsection{Instability of Actor-Critic Algorithms}
Our framework further provides tools for analyzing the instability that can arise in conventional actor-critic methods frequently used in RLHF, such as PPO. We align with the RLHF methodology, concentrating on the reinforcement learning fine-tuning phase as outlined in Section~\ref{section:prelims}. This analysis can be related to the control-as-inference paradigm \cite{levine2018reinforcement} for the constrained RL setting formulated in \ref{eq:RL}. Given a parameterized policy $\pi_{\theta}(y\mid x)$, the objective is to minimize $\mathbb{D}_{\text{KL}}[\pi_{\theta}(y|x) \mid \mid \pi^*(y\mid x)]$, where $\pi^*$ is the optimal policy derived in Eq.~\ref{eq:optimum_model} corresponding to the reward $r_{\phi}(y, x)$. Performing appropriate manipulations yields the following optimization objective:

\begin{equation}\label{eq:AC}
    \max_{\pi_{\theta}}\mathbb{E}_{\pi_{\theta}(y\mid x)}\bigg[\underbrace{r_{\phi}(x, y) -\beta\log\sum_{y}\pi_\text{ref}(y\mid x)\exp\left(\frac{1}{\beta}r_{\phi}(x, y)\right)}_{f(r_{\phi}, \pi_\text{ref}, \beta)} - \underbrace{\beta\log\frac{\pi_{\theta}(y\mid x)}{\pi_\text{ref}(y\mid x)}}_{\text{KL}}\bigg]
\end{equation}

This objective mirrors that employed in earlier research  
\citep{ziegler2020finetuning, stiennon2022learning, bai2022training, ouyang2022training}, where optimization is performed using the DPO-equivalent reward structure for the reward class $r_{\phi}$. Here, the normalization component within $f(r_{\phi}, \pi_\text{ref}, \beta)$ can be seen as analogous to the soft value function associated with the reference policy $\pi_\text{ref}$. Although this term does not alter the optimal policy, omitting it can result in increased variance in policy gradient estimates, thereby leading to instability during training. One way to manage the normalization factor is by incorporating a learned value function, though this itself poses significant optimization challenges. Alternatively, earlier studies have normalized rewards via a baseline derived from human completions, which serves as a one-sample Monte Carlo estimator for the normalization term. By contrast, DPO reparameterization leads to a reward structure that is inherently baseline-free.

\section{Experiments}
Here, we present empirical results evaluating DPO’s capability for learning policies directly from preference data. We first investigate, in a controlled text-generation benchmark, how DPO manages the tradeoff between reward maximization and KL-divergence regularization with respect to the reference policy, relative to typical preference optimization approaches such as PPO. Subsequently, we test DPO on more complex tasks involving larger models and challenging RLHF problems, including tasks in summarization and dialogue. Our findings indicate that, with minimal hyperparameter adjustment, DPO consistently matches or exceeds the performance of strong baselines like RLHF with PPO, and also outperforms methods selecting the best of $N$ trajectories using a learned reward model. Prior to discussing these experimental outcomes, we detail our experimental methodology, with supplementary information provided in Appendix~\ref{app:exp_details}.

\textbf{Tasks.} Our study investigates three distinct open-ended text generation tasks. Across all tasks, the various algorithms aim to learn a policy from a preference dataset $\mathcal{D}=\bigl\{x^{(i)}, y_w^{(i)}, y_l^{(i)}\bigr\}_{i=1}^N$. In the \textbf{controlled sentiment generation} task, $x$ is the introductory portion of a movie review sourced from the IMDb dataset \cite{maas-EtAl:2011:ACL-HLT2011}, and the objective for the policy is to continue this with $y$ exhibiting positive sentiment. To ensure controlled assessment, this setup employs \textit{generated} preference pairs, evaluated using a pre-trained sentiment classifier such that $p(\text{positive}\mid x,y_w)>p(\text{positive}\mid x,y_l)$. For SFT, we perform fine-tuning of GPT-2-large on the training portion of the IMDB dataset until the model reaches convergence (see further details in App~\ref{app:sentiment_details}). For the \textbf{summarization} task, $x$ refers to a Reddit forum post, with the policy required to generate a summary $y$ capturing the post’s main ideas. Following established methodologies, we utilize the Reddit TL;DR summarization dataset \citep{volske-etal-2017-tl} and leverage human preference annotations provided by \citeauthor{stiennon2022learning}. An SFT model is employed here, which has been fine-tuned on human-created forum summaries\footnote{\url{https://huggingface.co/CarperAI/openai_summarize_tldr_sft}}, utilizing the TRLX \citep{leandro_von_werra_2023_7790115} RLHF framework. The preferences for this dataset were collected by \citeauthor{stiennon2022learning} on outputs from a different SFT model, trained on comparable data. In the case of \textbf{single-turn dialogue}, $x$ is a user-posed query, varying from scientific inquiries to personal advice requests, where the model must produce an informative and engaging response $y$. The data source here is the Anthropic Helpful and Harmless dialogue dataset \citep{bai2022training}, which includes 170k conversations between a human and an AI assistant, each ending in two model-generated responses and an associated preference indicating which the human favored. No pre-trained SFT model is available for this scenario, so we construct the SFT model by further tuning a standard language model solely on preferred responses.

\textbf{Evaluation.} We employ two distinct evaluation strategies in our experiments. To assess how well each algorithm optimizes the constrained reward maximization objective within a controlled sentiment generation framework, we chart the set of attainable reward and KL-divergence trade-offs for each method. This analysis is feasible in our setting because we possess access to the exact reward signal (i.e., the sentiment classifier). In contrast, in practical, real-world scenarios, where the true reward function remains unknown, we instead assess algorithms by their \textit{win rate} relative to a baseline policy. Here, GPT-4 acts as a surrogate for human evaluators, rating summary quality in the summarization task and response helpfulness in the single-turn dialogue task. For the summarization experiments, we designate the test set’s reference summaries as the baseline, whereas in the dialogue experiments, the ground-truth preferred responses from the test set are used as baselines. Although recent literature points to large language models surpassing standard automatic metrics as evaluators \citep{Chen2023ExploringTU}, we further support our choice to use GPT-4 for evaluation with a dedicated human study (see Sec.~\ref{sec:human-judgments}). Our analysis indicates that GPT-4’s judgments align closely with those of human annotators, often matching or exceeding the consistency found between different human raters.

\textbf{Methods.} Alongside DPO, we benchmark a variety of established techniques for training language models aligned with human preferences. For baseline prompting, we apply zero-shot prompting to \textbf{GPT-J} \citep{gpt-j} in the summarization domain and use two-shot prompting with \textbf{Pythia-2.8B} \citep{biderman2023pythia} in dialogue generation. We further evaluate the \textbf{SFT} model and introduce \textbf{Preferred-FT}, which is obtained by supervised fine-tuning on selected outputs $y_w$—drawing from either the SFT model (in the cases of sentiment control and summarization) or a generic LM (in the single-turn dialogue setting). Another approach considered is the pseudo-supervised \textbf{Unlikelihood}~\citep{welleck2019neural} method, which aims to increase the model’s probability on $y_w$ while explicitly decreasing probability assigned to $y_l$, employing an adjustable scaling factor $\alpha\in[0,1]$ for the unlikelihood loss component. We also evaluate \textbf{PPO} \citep{schulman2017proximal}, where the reward signal is derived from learned preferences, and \textbf{PPO-GT}, an oracle variant leveraging the true reward in the sentiment-controlled setup. For sentiment experiments, two PPO-GT variants are considered: a standard off-the-shelf implementation \cite{leandro_von_werra_2023_7790115} and an enhanced version with normalized rewards and optimized hyperparameters to boost efficacy—such hyperparameter adjustments are also employed when training standard PPO with learned rewards. As an additional point of comparison, we utilize the \textbf{Best of $N$} baseline, wherein $N$ candidate outputs are generated from the SFT (or Preferred-FT in dialogue) model and the one scoring highest under the learned reward function is selected. While this approach effectively decouples reward model quality from PPO-style optimization, it is prohibitively expensive for large $N$, as it necessitates sampling $N$ completions per input at evaluation time.

\subsection{How well can DPO optimize the RLHF objective?}

The reward maximization objective in RLHF algorithms typically incorporates a KL constraint to ensure that the learned policy remains sufficiently close to the reference policy, trading off between maximizing reward and limiting policy divergence. Thus, when evaluating and comparing different algorithms, it is critical to consider not only the magnitude of the reward attained but also the associated KL divergence; in many cases, achieving higher rewards at the cost of greatly increased KL is suboptimal. In Figure~\ref{fig:frontier-tldr-main}, we illustrate the relationship between reward and KL for various algorithms in the sentiment classification scenario. For each method, we conduct several training runs, each with a distinct setting of the policy conservativeness hyperparameter (target KL $\in\{3,6,9,12\}$ for PPO, $\beta \in \{0.05,0.1,1,5\}$, $\alpha\in\{0.05,0.1,0.5,1\}$ for unlikelihood, and differing random seeds for preferred-FT). This results in a total of 22 experimental runs. After every 100 training updates up to convergence, we evaluate each trained policy over a test prompt set, measuring the mean reward as defined by the ground-truth reward function, and calculating the average sequence-level KL\footnote{Here, this refers to the aggregate sum of the KL-divergence across all timesteps in a sequence.} relative to the reference policy $\text{KL}\left(\pi\mid \mid \pi_\text{ref}\right)$. The experiments reveal that DPO establishes a highly efficient reward-KL frontier, delivering superior reward outcomes while retaining low KL values. This is a particularly important finding for several reasons. Firstly, although both DPO and PPO target the same optimization objective, DPO exhibits markedly better efficiency, as its reward/KL performance consistently surpasses that of PPO. Secondly, DPO demonstrates a superior reward-KL frontier even when PPO has access to exact, ground-truth rewards (PPO-GT).

\subsection{Can DPO scale to real preference datasets?}
\label{sec:dpo-real-datasets}
We subsequently assess DPO's effectiveness when fine-tuned on real-world datasets involving summarization and single-turn dialogue tasks. In summarization, established automatic metrics such as ROUGE have been shown to align poorly with actual human judgments~\citep{stiennon2022learning}, and past research demonstrated that leveraging PPO-based fine-tuning on human preference data can enhance summary quality. To benchmark methods, we generate outputs on the test split of the TL;DR summarization dataset and measure the mean win rate relative to the dataset's reference completions. For each method, outputs are produced using sampling temperatures ranging between 0.0 and 1.0, with corresponding win rates depicted in Figure~\ref{fig:frontier-tldr-main} (right). DPO, PPO, and Preferred-FT are all fine-tuned from the same GPT-J SFT model\footnote{\url{https://huggingface.co/CarperAI/openai_summarize_tldr_sft}}. Our results indicate that DPO achieves an approximate 61\% win rate at a temperature of 0.0, outperforming PPO, which attains around 57\% at its best sampling temperature, also 0.0. Furthermore, DPO secures a greater peak win rate compared to the best-of-$N$ baseline approach. It should be noted that no significant optimization was performed for DPO’s $\beta$ hyperparameter; thus, the reported outcomes might not represent the method’s full capability. Additionally, DPO demonstrates a higher degree of stability across varying sampling temperatures, whereas PPO’s performance can drop to the level of the base GPT-J model under higher temperatures. Preferred-FT does not yield noticeable improvements over the initial SFT baseline. In direct comparisons through human evaluation, discussed further in Section~\ref{sec:human-judgments}, completions generated by DPO at a sampling temperature of 0.25 were favored 58\% of the time over PPO outputs sampled at the same temperature.

For the single-turn dialogue experiments, we conduct assessments on the subset of the Anthropic HH dataset \citep{bai2022training} test split that involves a single exchange between human and assistant. In our GPT-4 assessments, we employ the preferred responses from the test set as ground truth references for calculating the win rates of the evaluated methods. In the absence of a commonly established SFT baseline for this context, we begin with the pre-trained Pythia-2.8B model, apply Preferred-FT to adapt a reference model using the selected completions—ensuring the resulting outputs remain within the training distribution—and subsequently train the model utilizing DPO. Additionally, we benchmark against the strongest performer among 128 Preferred-FT generated completions (as our investigation indicates the Best of $N$ baseline’s performance levels off at $N=128$ for this dataset; refer to Appendix Figure~\ref{fig:best-of-n}) and a two-shot prompted version of the original Pythia-2.8B model. Our findings show that, across optimal temperature selections, DPO matches or exceeds the results achieved by the best alternative approaches. We further analyze an RLHF model that employs PPO and has been trained on the Anthropic HH dataset\footnote{\url{https://huggingface.co/reciprocate/ppo_hh_pythia-6B}} from a widely used repository\footnote{\url{https://github.com/CarperAI/trlx/tree/main/examples/hh}}, but we do not identify any prompt or temperature configurations that result in improved performance over the base Pythia-2.8B model. Drawing on insights from TL;DR as well as the observation that both methods are tuned to the same reward function, we use the Best of 128 as an approximate indicator of PPO-level efficacy. In summary, DPO stands out as the only computationally efficient strategy that surpasses the reference completions on the Anthropic HH dataset, achieving results on par with or exceeding the much more resource-intensive Best of 128 baseline. As further illustrated in Figure~\ref{fig:dialogue-main}, DPO rapidly approaches its optimal performance.

\subsection{Generalization to a new input distribution}

To more thoroughly assess how PPO and DPO handle shifts in input distribution, we apply the policies trained in our Reddit TL;DR summarization setup to a distinct data domain: the test set of the CNN/DailyMail news articles dataset \citep{nallapati-etal-2016-abstractive}. We use the optimal sampling temperatures identified for TL;DR (0 and 0.25) for this evaluation. Table~\ref{tab:ood} summarizes these findings. To measure performance, we calculate the rate at which GPT-4 prefers model-generated summaries over ground-truth summaries using a variant of the GPT-4 (C) prompt, modified to substitute “news article” for “forum post.” On this new domain, DPO maintains a sizable performance edge over PPO. These findings provide early evidence that DPO can match or surpass PPO in out-of-distribution generalization, even though DPO is trained without the extra unlabeled prompts from Reddit TL;DR that PPO leverages.

\subsection{Validating GPT-4 judgments with human judgments}
\label{sec:human-judgments}
To determine the trustworthiness of GPT-4 as an evaluator, we conduct a human-subject experiment utilizing outcomes from the TL;DR summarization study, and compare GPT-4’s ratings across two distinct prompts. The \textbf{GPT-4 (S)} (simple) prompt requests a decision about which summary best encapsulates the essential content of the post. In contrast, the \textbf{GPT-4 (C)} (concise) prompt adds an explicit preference for conciseness; we analyze this prompt because the simple version causes GPT-4 to favor lengthier, more redundant summaries than those preferred by humans. Full prompt texts appear in Appendix~\ref{app:prompts}. We perform three matchups: the highest-performing method (DPO, temp. 0.25), the lowest (PPO, temp. 1.0), and an intermediate approach (SFT, temp. 0.25), each compared against PPO sampled greedily (its strongest setting). These pairings were selected to span a broad range of summary quality; for the DPO and PPO-1 comparisons, multiple human ratings per example were gathered, though human ratings are otherwise sparse. Across both prompts, agreement rates between GPT-4 and human evaluators are comparable to agreement rates among different human annotators, indicating GPT-4 can reliably approximate human preferences. Furthermore, win rates from the \textbf{GPT-4 (C)} prompt consistently align better with human judgments; thus, this prompt is adopted for primary results in Section~\ref{sec:dpo-real-datasets}. Further methodological details, including the rater-facing web interface and roster of participating volunteers, can be found in Appendix~\ref{app:human-study}.

\section{Discussion}
Preference-based learning offers an effective and scalable approach for developing capable and value-aligned language models. In this work, we have presented DPO, a straightforward methodology that facilitates training language models from preferences without the need for reinforcement learning. Instead of transforming the preference learning task into a conventional RL framework to leverage pre-existing RL techniques, DPO establishes a correspondence between language model policies and reward functions. This allows direct alignment of language model behavior with human preferences using a simple cross-entropy objective, thus removing the requirement for reinforcement learning and maintaining generality. Remarkably, DPO achieves performance on par with or exceeding leading RLHF methods—such as those utilizing PPO—without the need for extensive hyperparameter adjustment, thus significantly lowering the practical barriers to preference-based language model training.

\textbf{Limitations \& Future Directions.} Several open questions emerge from our findings and provide avenues for further investigation. One such question concerns the generalization capacity of DPO-trained policies to out-of-distribution inputs, especially compared to policies trained with explicit reward signals; preliminary evidence suggests DPO generalizes comparably to PPO-trained policies, but a thorough exploration is warranted. Another important line of inquiry is whether leveraging self-labeling based on the DPO policy can effectively exploit unlabeled data. Furthermore, the dynamics of reward over-optimization under direct preference optimization require examination—specifically, whether the minor performance reduction observed in Figure~\ref{fig:dialogue-main}-right may be attributed to this effect. Although our experiments are confined to models with up to 6B parameters, scaling DPO to much larger, state-of-the-art architectures remains a promising challenge. Regarding evaluation protocols, we observe that win rates as assessed by GPT-4 are sensitive to prompt formulation; determining strategies for eliciting consistent and reliable assessments from automated evaluators presents a fertile area for subsequent study. Lastly, DPO’s applicability extends beyond preference-driven language modeling, with potential utility in the training of generative models across a variety of domains. 

\end{document}